% Generated by scripts/assemble.py; corrections belong in source/corrections.json.
\examhead{3}{北京师范大学 2026 年数学分析试卷}{10}
% p011-b12; source PDF page 11
\par\Needspace{5\baselineskip}
\par\Needspace{6\baselineskip}
\blockanchor{p011-b12}
\majorquestion{一}(15 分) 证明:\par
\nopagebreak[4]
% p011-b13; source PDF page 11
\blockanchor{p011-b13}
\displayformula{1 - \frac {1}{2} + \frac {1}{3} - \frac {1}{4} + \dots + \frac {1}{2025} - \frac {1}{2026} <   \frac {\sqrt {2}}{2}.}
% p011-b14; source PDF page 11
\par\Needspace{5\baselineskip}
\par\Needspace{6\baselineskip}
\blockanchor{p011-b14}
\majorquestion{二}(15 分) 设 $f: \mathbb{R} \times \mathbb{R} \to \mathbb{R}$ 满足对任意的 $a \in \mathbb{R}$, 有 $f(0, a) \leq f(0, 0) \leq f(a, 0)$. 证明:\par
\nopagebreak[4]
% p011-b15; source PDF page 11
\blockanchor{p011-b15}
\displayformula{\sup _ {b \in \mathbb {R}} \inf _ {a \in \mathbb {R}} f (a, b) = \inf _ {a \in \mathbb {R}} \sup _ {b \in \mathbb {R}} f (a, b).}
% p011-b16; source PDF page 11
\par\Needspace{5\baselineskip}
\blockanchor{p011-b16}
\majorquestion{三}(15 分) 求极限 $\lim_{x\to +\infty}x^2\left[\left(\frac{3x^3 - 7}{3x^3 - x}\right)^{\frac{1}{5}} - x\sin \frac{1}{x}\right].$\par
% p011-b17; source PDF page 11
\par\Needspace{5\baselineskip}
\blockanchor{p011-b17}
\majorquestion{四}(15 分) 已知 $f(x) = \sum_{n=1}^{\infty} \frac{\cos(nx)}{\sqrt{n^3 + n}}$. 证明:\par
% p011-b18; source PDF page 11
\blockanchor{p011-b18}
(1) (7 分) $f(x)$ 在 $(- \infty, + \infty)$ 上连续.\par
% p011-b19; source PDF page 11
\blockanchor{p011-b19}
(2) (8 分) $\frac{\sqrt{2}}{2} - \frac{1}{15} < \int_{0}^{\frac{\pi}{2}} f(x) \, \mathrm{d}x < \frac{\sqrt{2}}{2}$.\par
% p011-b20; source PDF page 11
\par\Needspace{5\baselineskip}
\par\Needspace{6\baselineskip}
\blockanchor{p011-b20}
\majorquestion{五}(15 分) 设 $f(x)$ 在 $[1, +\infty)$ 上非负单调递减, 证明: $\lim_{n \to \infty} \left[ \sum_{k=1}^{n} f(k) - \int_{1}^{n} f(x) \, \mathrm{d}x \right]$ 收敛.
\majorquestion{六}(15 分) 计算广义积分 $\int_{1}^{+\infty} \frac{\{x\}}{x^3} \, \mathrm{d}x$, 其中 $\{x\}$ 表示 $x$ 的小数部分, 例如 $\{\pi\} = \pi - 3$.
\majorquestion{七}(15 分) 设 $\delta \in [0, +\infty)$, 讨论函数\par
\nopagebreak[4]
% p011-b21; source PDF page 11
\blockanchor{p011-b21}
\displayformula{f (x, y) = \left\{ \begin{array}{l l} \frac {\sin (x ^ {1 + \delta} y)}{x ^ {2} + y ^ {2}}, & (x, y) \neq (0, 0); \\ 0, & (x, y) = (0, 0). \end{array} \right.}
% p011-b22; source PDF page 11
\blockanchor{p011-b22}
的连续性, 可微性.\par
% p012-b00; source PDF page 12
\par\Needspace{5\baselineskip}
\par\Needspace{6\baselineskip}
\blockanchor{p012-b00}
\majorquestion{八}(15 分) 设 $f(x)$ 在 $[0, +\infty)$ 上连续.\par
\nopagebreak[4]
% p012-b01; source PDF page 12
\blockanchor{p012-b01}
\displayformula{f (t) = t + \frac {1}{4 \pi} \iiint_ {x ^ {2} + y ^ {2} + z ^ {2} \leq t ^ {2}} \frac {f (\sqrt {x ^ {2} + y ^ {2} + z ^ {2}})}{x ^ {2} + y ^ {2} + z ^ {2}} d x d y d z, t \in [ 0, + \infty).}
% p012-b02; source PDF page 12
\blockanchor{p012-b02}
求 $f(t)$ .\par
% p012-b03; source PDF page 12
\par\Needspace{5\baselineskip}
\par\Needspace{6\baselineskip}
\blockanchor{p012-b03}
\majorquestion{九}(15 分) 求曲线积分\par
\nopagebreak[4]
% p012-b04; source PDF page 12
\blockanchor{p012-b04}
\displayformula{I = \int_ {L} \frac {(x - y) \mathrm{d} x + (x + 4 y) \mathrm{d} y}{x ^ {2} + 4 y ^ {2}}.}
% p012-b05; source PDF page 12
\blockanchor{p012-b05}
其中 L 从  $(1,0)$  沿上半圆周  $x^{2}+y^{2}=1$  到  $(-1,0)$ .\par
% p012-b06; source PDF page 12
\par\Needspace{5\baselineskip}
\blockanchor{p012-b06}
\majorquestion{十}(15 分) 在空间直角坐标系中, 已知方程 $\cos(3x + 2z) = h(4y - z)$, 且函数 $h$ 在 $\mathbb{R}$ 上可微, 证明: 上述方程表示空间中的一个柱面.\par
